% GENERATED FILE. Edit canonical lessons, then run npm run generate:books.
% canonical-source-hash: fnv1a64:45e7a94b5add5381
% canonical-lessons: GU-C191-dhirvu, GU-C191-vyaj, GU-C191-bacat, GU-C191-napho, GU-C191-khot

\chapter{Lending at interest}
\label{ch:gu-lending-at-interest}
\hlchaptermodality{191}

\begin{chapteropening}
\textbf{By the end of this chapter:} \emph{I can talk about lending, interest, savings, profit and loss.}

Complete the last lesson of chapter 191: I can talk about lending, interest, savings, profit and loss.
\end{chapteropening}

\section[\texorpdfstring{\gu{ધીરવું} (dhīrvũ)}{ધીરવું (dhīrvũ)}]{\gu{ધીરવું} (dhīrvũ) — to lend}

\label{lesson:GU-C191-dhirvu}



\begin{quote}
Before the new one: say the Gujarati for to cash a cheque, then the Gujarati for to borrow.
\end{quote}

\subsection*{You'll want to know: \gu{ધીરવું}}
\textbf{\gu{ધીરવું}} — \emph{dhīrvũ} — \textquotedblleft{}to lend\textquotedblright{}.

\begin{quote}
\textbf{\gu{તે} \gu{મિત્રને} \gu{પૈસા} \gu{ધીરે} \gu{છે}.} — \emph{te mitrane paisā dhīre che} — \textquotedblleft{}He lends money to a friend.\textquotedblright{}
\end{quote}

\begin{cousinweb}
\textbf{\gu{ધીરવું}} (\emph{dhīrvũ}) is a tadbhava. CDIAL lists it, glossed \textquotedblleft{}to trust, lend\textquotedblright{}, under Sanskrit \emph{dh\'{\={\i}}ra}, \textquotedblleft{}wise; steady, resolute\textquotedblright{} (CDIAL 6817).
\end{cousinweb}

\begin{grammarlens}[title={What you've built}]
One more word for everyday life at A2.
\end{grammarlens}

\subsection*{Your turn}
\begin{itemize}
\raggedright
  \item \emph{Say it:} \emph{dhīrvũ}
  \item \emph{Say it:} \emph{dhīrvũ}, once more
  \item \emph{Say it:} \emph{te mitrane paisā dhīre che} — He lends money to a friend.
  \item \emph{Recall:} say the Gujarati for to cash a cheque, then the Gujarati for to borrow, then say \emph{dhīrvũ} again
\end{itemize}

\begin{tcolorbox}[breakable,colback=teal!4,colframe=teal!35!black,title={Before you move on}]
What does \gu{ધીરવું} mean? (\textquotedblleft{}To lend\textquotedblright{}.) Say it once more.
\end{tcolorbox}
\section[\texorpdfstring{\gu{વ્યાજ} (vyāj)}{વ્યાજ (vyāj)}]{\gu{વ્યાજ} (vyāj) — interest on money}

\label{lesson:GU-C191-vyaj}



\begin{quote}
Before the new one: say the Gujarati for to borrow, then the Gujarati for to lend.
\end{quote}

\subsection*{You'll want to know: \gu{વ્યાજ}}
\textbf{\gu{વ્યાજ}} — \emph{vyāj} — \textquotedblleft{}interest on money\textquotedblright{}.

\begin{quote}
\textbf{\gu{આ} \gu{ખાતામાં} \gu{વ્યાજ} \gu{ઓછું} \gu{છે}.} — \emph{ā khātāmā̃ vyāj ochũ che} — \textquotedblleft{}This account pays little interest.\textquotedblright{}
\end{quote}

\begin{cousinweb}
The origin of \textbf{\gu{વ્યાજ}} (\emph{vyāj}) is not settled. It matches Sanskrit \emph{vyāja}, \textquotedblleft{}pretext\textquotedblright{}, in form, but CDIAL has no entry for \emph{vyāja} and no Gujarati word for interest, so the link is unconfirmed.
\end{cousinweb}

\begin{grammarlens}[title={What you've built}]
One more word for everyday life at A2.
\end{grammarlens}

\subsection*{Your turn}
\begin{itemize}
\raggedright
  \item \emph{Say it:} \emph{vyāj}
  \item \emph{Say it:} \emph{vyāj}, once more
  \item \emph{Say it:} \emph{ā khātāmā̃ vyāj ochũ che} — This account pays little interest.
  \item \emph{Recall:} say the Gujarati for to borrow, then the Gujarati for to lend, then say \emph{vyāj} again
\end{itemize}

\begin{tcolorbox}[breakable,colback=teal!4,colframe=teal!35!black,title={Before you move on}]
What does \gu{વ્યાજ} mean? (\textquotedblleft{}Interest on money\textquotedblright{}.) Say it once more.
\end{tcolorbox}
\section[\texorpdfstring{\gu{બચત} (bacat)}{બચત (bacat)}]{\gu{બચત} (bacat) — savings}

\label{lesson:GU-C191-bacat}



\begin{quote}
Before the new one: say the Gujarati for to lend, then the Gujarati for interest on money.
\end{quote}

\subsection*{You'll want to know: \gu{બચત}}
\textbf{\gu{બચત}} — \emph{bacat} — \textquotedblleft{}savings\textquotedblright{}.

\begin{quote}
\textbf{\gu{મારી} \gu{બચત} \gu{બેંકમાં} \gu{છે}.} — \emph{mārī bacat bẽkmā̃ che} — \textquotedblleft{}My savings are in the bank.\textquotedblright{}
\end{quote}

\begin{cousinweb}
\textbf{\gu{બચત}} (\emph{bacat}) is built on the verb \emph{bacvũ}, \textquotedblleft{}to be left over\textquotedblright{}, with the noun ending \emph{-at}: what is left over. \emph{bacvũ} came into Gujarati from Hindi and goes back to Sanskrit \emph{váñcati}, \textquotedblleft{}escapes from\textquotedblright{} (CDIAL 11208).
\end{cousinweb}

\begin{grammarlens}[title={What you've built}]
One more word for everyday life at A2.
\end{grammarlens}

\subsection*{Your turn}
\begin{itemize}
\raggedright
  \item \emph{Say it:} \emph{bacat}
  \item \emph{Say it:} \emph{bacat}, once more
  \item \emph{Say it:} \emph{mārī bacat bẽkmā̃ che} — My savings are in the bank.
  \item \emph{Recall:} say the Gujarati for to lend, then the Gujarati for interest on money, then say \emph{bacat} again
\end{itemize}

\begin{tcolorbox}[breakable,colback=teal!4,colframe=teal!35!black,title={Before you move on}]
What does \gu{બચત} mean? (\textquotedblleft{}Savings\textquotedblright{}.) Say it once more.
\end{tcolorbox}
\section[\texorpdfstring{\gu{નફો} (napho)}{નફો (napho)}]{\gu{નફો} (napho) — profit}

\label{lesson:GU-C191-napho}



\begin{quote}
Before the new one: say the Gujarati for interest on money, then the Gujarati for savings.
\end{quote}

\subsection*{You'll want to know: \gu{નફો}}
\textbf{\gu{નફો}} — \emph{napho} — \textquotedblleft{}profit\textquotedblright{}.

\begin{quote}
\textbf{\gu{આ} \gu{સોદામાં} \gu{નફો} \gu{થયો}.} — \emph{ā sodāmā̃ napho thayo} — \textquotedblleft{}This deal made a profit.\textquotedblright{}
\end{quote}

\begin{cousinweb}
\textbf{\gu{નફો}} (\emph{napho}) is a loan from Arabic \emph{nafʿ}, \textquotedblleft{}profit, benefit\textquotedblright{}, which reached Gujarati through Persian. Its Arabic root is \emph{n-f-ʿ}, \textquotedblleft{}to be of use\textquotedblright{} (Lane).
\end{cousinweb}

\begin{grammarlens}[title={What you've built}]
One more word for everyday life at A2.
\end{grammarlens}

\subsection*{Your turn}
\begin{itemize}
\raggedright
  \item \emph{Say it:} \emph{napho}
  \item \emph{Say it:} \emph{napho}, once more
  \item \emph{Say it:} \emph{ā sodāmā̃ napho thayo} — This deal made a profit.
  \item \emph{Recall:} say the Gujarati for interest on money, then the Gujarati for savings, then say \emph{napho} again
\end{itemize}

\begin{tcolorbox}[breakable,colback=teal!4,colframe=teal!35!black,title={Before you move on}]
What does \gu{નફો} mean? (\textquotedblleft{}Profit\textquotedblright{}.) Say it once more.
\end{tcolorbox}
\section[\texorpdfstring{\gu{ખોટ} (khoṭ)}{ખોટ (khoṭ)}]{\gu{ખોટ} (khoṭ) — loss, shortfall}

\label{lesson:GU-C191-khot}



\begin{quote}
Before the new one: say the Gujarati for savings, then the Gujarati for profit.
\end{quote}

\subsection*{You'll want to know: \gu{ખોટ}}
\textbf{\gu{ખોટ}} — \emph{khoṭ} — \textquotedblleft{}loss, shortfall\textquotedblright{}.

\begin{quote}
\textbf{\gu{વેપારીને} \gu{ખોટ} \gu{ગઈ}.} — \emph{vepārīne khoṭ gaī} — \textquotedblleft{}The merchant made a loss.\textquotedblright{}
\end{quote}

A relative of \textbf{\gu{ખોટું}}, \emph{khoṭũ}, 'false'.

\begin{cousinweb}
\textbf{\gu{ખોટ}} (\emph{khoṭ}) is a tadbhava. CDIAL lists it, glossed \textquotedblleft{}mistake, loss, want\textquotedblright{}, under \emph{*khōṭi} and \emph{*khōṭṭa}, \textquotedblleft{}blemish\textquotedblright{}, reconstructed forms, beside \emph{khoṭũ}, \textquotedblleft{}alloyed, bad\textquotedblright{} (CDIAL 3931).
\end{cousinweb}

\begin{grammarlens}[title={What you've built}]
One more word for everyday life at A2.
\end{grammarlens}

\subsection*{Your turn}
\begin{itemize}
\raggedright
  \item \emph{Say it:} \emph{khoṭ}
  \item \emph{Say it:} \emph{khoṭ}, once more
  \item \emph{Say it:} \emph{vepārīne khoṭ gaī} — The merchant made a loss.
  \item \emph{Recall:} say the Gujarati for savings, then the Gujarati for profit, then say \emph{khoṭ} again
\end{itemize}

\begin{tcolorbox}[breakable,colback=teal!4,colframe=teal!35!black,title={Before you move on}]
What does \gu{ખોટ} mean? (\textquotedblleft{}Loss, shortfall\textquotedblright{}.) Say it once more.
\end{tcolorbox}
